\section{稳定匹配与延迟接受}
\label{sec:04-Discrete-Mathematics/03-Graph-Theory/06}

\subsection{三个人、三只猫、一个困境}

你运营一家领养站。今天来了三个人，三只猫：

\begin{center}
\begin{tabular}{c|c}
人 & 对猫的偏好（从左到右是最喜欢到最不喜欢）\\
\hline
人 1 & 猫 A > 猫 B > 猫 C\\
人 2 & 猫 A > 猫 C > 猫 B\\
人 3 & 猫 B > 猫 A > 猫 C
\end{tabular}
\quad
\begin{tabular}{c|c}
猫 & 对人的偏好\\
\hline
猫 A & 人 2 > 人 1 > 人 3\\
猫 B & 人 1 > 人 3 > 人 2\\
猫 C & 人 3 > 人 2 > 人 1
\end{tabular}
\end{center}

任务很简单：配出三个``人--猫''对，每个人都分到一只猫，每只猫都分到一个人。

这事看起来毫无难度——你就从第一个人开始挨个挑，先到先得，猫被领走就顺延下一个，一分钟就配完了。真的配完了吗？

\subsection{贪婪先到先得：快，但不稳}

按编号顺序，人 1 先挑。他最喜欢猫 A，猫 A 还在——到手。人 2 也最喜欢猫 A，但猫 A 刚被领走，他退而求猫 C。人 3 喜欢猫 B，猫 B 还未被领——顺利配对。

结果：(人 1, 猫 A), (人 2, 猫 C), (人 3, 猫 B)。耗时十秒，匹配完成。

现在从猫的视角审视这份名单。猫 A 被分给了人 1，但它自己的偏好是 人 2 > 人 1 > 人 3——它更喜欢人 2。再看人 2：他分到了猫 C，但他也更喜欢猫 A 胜过猫 C。

人 2 和猫 A 互相对视了一眼。双方都更愿意跟对方配对，胜过各自当前的对象。

这不是心理活动——这是结构性的裂隙。只要有一对未配在一起的成员同时更偏爱对方，他们就有动力私下交换，当前的配对方案会自行解体。这样的对子叫阻塞对（blocking pair）：猫 A 和它的分到对象人 1 之间没有阻塞，但猫 A 和没分给它的人 2 之间有阻塞。一个存在阻塞对的匹配被称为不稳定的。

贪婪的根伤在哪？不是算错了偏好，也不是程序 bug。它失败是因为一个更根本的原因：最早挑的人占了当时最好的，但他未必是对方眼中的最佳人选。早到的优势被锁死，后来者只能从剩余选项里挑，而剩余选项的主人们可能眼睁睁看着自己更喜欢的人被更早的对手领走、无力反抗。时序上的先后成了不可逆的权力——这恰恰是``稳定''一词要挡住的漏洞。

\subsection{什么才算稳定：一个结构性的最低门槛}

先给一个准确的定义。一个匹配中没有阻塞对，就叫稳定匹配。

注意这个定义不关心三件事。不关心总满意度——可能存在另一个匹配让所有人的偏好排名之和更小，但只要其中某对成员有动力脱队，它就不稳定。不关心公平——同为人，张三可能分到自己的第一志愿，李四只能分到第三志愿，但只要李四想撬的人根本看不上他，这仍然是稳定匹配。不关心唯一性——同一个问题可能有好几个稳定匹配，彼此的总满意度可以差很多。

稳定只排除一种局部偏离：双方都有动力私下改配。它是一个结构性的最低门槛。一个匹配如果连这个都做不到，推向市场立刻散架。但它也只是最低门槛——跨过它，离``好''还很远。

\subsection{换个思路：不急着定案}

贪婪的问题出在``一配对就锁死''。猫 A 一旦跟了人 1，人 2 再有更好的条件也撬不动。但如果暂时不锁呢？

想象一种完全不同的流程。每个人向自己喜欢的猫发出提议，猫不立刻答应——猫把当前收到的所有提议中自己最喜欢的那个暂时保留下来，其他的拒绝掉。被拒绝的人划掉这个选项，转向次优选择继续提议。猫手里的暂定人选可以在更好的提议到来时被替换。

这个直觉可以立刻跑一遍上面的例子。

\subsubsection{第一轮}

人 1 向猫 A 提议。猫 A 手头空着，暂时保留人 1。
人 2 向猫 A 提议。猫 A 比较人 1 和人 2——按它的偏好，人 2 > 人 1。猫 A 改保留人 2，拒绝人 1。
人 3 向猫 B 提议。猫 B 手头空着，暂时保留人 3。

\subsubsection{第二轮}

人 1 被猫 A 拒绝。他从偏好列表上划掉猫 A，下一个最喜欢的是猫 B。人 1 向猫 B 提议。
猫 B 比较手头的人 3 和新人 1——按它的偏好，人 1 > 人 3。猫 B 改保留人 1，拒绝人 3。
人 2 和猫 A 暂配，本轮无事。

\subsubsection{第三轮}

人 3 被猫 B 拒绝。划掉猫 B，下一个是猫 A。人 3 向猫 A 提议。
猫 A 比较手头的人 2 和新人 3——按偏好，人 2 > 人 3。猫 A 继续保留人 2，拒绝人 3。

\subsubsection{第四轮}

人 3 被猫 A 拒绝。划掉猫 A，下一个是猫 C。人 3 向猫 C 提议。
猫 C 还没收过任何提议，保留人 3。

\subsubsection{终止}

现在人 1 与猫 B 互配，人 2 与猫 A 互配，人 3 与猫 C 互配。无人剩下可提议对象，算法终止。

验证稳定性：猫 A 更偏爱人 2 胜过人 1 和人 3——它已经配给了人 2，不可能被撬动。猫 B 更偏爱人 1 胜过人 3 和人 2——它已经配给了第一志愿。猫 C 更偏爱人 3——当前配对也是它的最优。人 2 和猫 A 互为第一志愿，人 1 得到猫 B 是第一志愿，人 3 得到猫 C 是第二志愿（猫 B 拒绝了他）。不存在任何阻塞对。

回头对比两个方案。贪婪给了：(人 1, 猫 A), (人 2, 猫 C), (人 3, 猫 B)——存在阻塞对 (人 2, 猫 A)，不匹配。延迟接受给了：(人 1, 猫 B), (人 2, 猫 A), (人 3, 猫 C)——稳定且也相当好。

\subsection{算法骨架与不变量的力量}

把这个直觉写成算法。一侧主动提议，另一侧被动选择。以人主动为例：

\begin{enumerate}
\tightlist
\item 每个尚未匹配且偏好列表还有未试对象的人，向列表中尚未拒绝过自己的第一个（即剩余中最喜欢的）对象提议；
\item 每个接收方暂时保留当前所有提议中最喜欢的一个，拒绝其余；
\item 若无人能继续提议，停止；否则回到第 1 步。
\end{enumerate}

终止是必然的。每个人的偏好列表有穷，每提一次议就划掉一个对象，最多提议 \(n^2\) 次。实际上通常远少于此。

Gale--Shapley 的妙处在两个方向相反的不变量。

提议方从偏好列表的高位向低位走。他被拒一次，只能往下走一步——他只能越找越差。接收方从低位向高位换。她手上的暂定人选只会被偏好列表中更高的人替换——她只会越换越好。

这两个单调性合在一起，掐死了阻塞对的存在空间。假设最终结果中存在阻塞对 \((a, b)\)——\(a\) 更喜欢 \(b\) 胜过自己的最终对象，同时 \(b\) 也更喜欢 \(a\) 胜过自己的最终对象。按算法，\(a\) 按偏好从高到低提议，他一定在向当前对象提议之前就已经向 \(b\) 提过议了。\(b\) 当时拒绝了 \(a\)，说明 \(b\) 那时手头有比 \(a\) 更喜欢的暂定对象。此后 \(b\) 的暂定对象只可能是偏好列表上更靠前的人——最终 \(b\) 的对象至少不比当初拒绝 \(a\) 时手中的差。但阻塞对要求 \(b\) 更喜欢 \(a\) 胜过最终对象，这与``越换越好''冲突。矛盾。

证明力量来自一组极简的约束：提议方单调下降，接收方单调上升，两条轨迹永远不会错肩出让阻塞对出现的空间。

\subsection{换一边提议，天翻地覆}

上面用的是``人主动''。如果换成猫主动：

\begin{itemize}
\tightlist
\item 猫 A 偏好 人 2 > 人 1 > 人 3——向人 2 提议，人 2 暂留。
\item 猫 B 偏好 人 1 > 人 3 > 人 2——向人 1 提议，人 1 暂留。
\item 猫 C 偏好 人 3 > 人 2 > 人 1——向人 3 提议，人 3 暂留。
\end{itemize}

一轮收工：每个人的第一志愿猫都向他们提议了，全部配对。结果：(猫 A, 人 2), (猫 B, 人 1), (猫 C, 人 3)。这也是稳定匹配，而且和上一轮的结果完全相同。

但在别的偏好转入下，换边会产生迥异的匹配。更一般地，这个不对称性有精确的含义：在严格完整的偏好设定下，主动提议侧得到的是所有稳定匹配中对本侧最优的结果——每一个提议者在任何其他稳定匹配中都不会得到比算法结果更喜欢的对象。被动接收侧得到的则是所有稳定匹配中对本侧最差的结果。

这意味着``产生稳定解''不等于机制中立。谁有提议权，整个解集就偏向哪一侧。这一点背后的博弈意味很重，我们搁到边界里讨论。

\subsection{承认边界：稳定只是最低门槛}

稳定匹配理论给了我们一个存在性定理（Gale--Shapley 总产出至少一个稳定匹配）和一个构造算法。但它的边界相当多，每一条都值得清醒地跨过去。

\textbf{偏好可能不是严格序。}真实的偏好经常含平局——两只猫在某人心中不分上下，或者一头猫在两个人之间完全无偏好。平局下``最优稳定匹配''和``最差稳定匹配''的优劣差距可能极大，且稳定匹配的结构比严格序复杂得多，不再呈现明确的格子结构。

\textbf{容量可能不是一对一的。}医院招住院医，一家医院可以收多名医学生。这时算法只需让接收方以容量为限，暂留多个对象而非一个，多余容量按偏好排名取优。核心不变量仍然成立：接收方手中的``最差暂定对象''在偏好上单调上升。

\textbf{有人可能谎报偏好。}这是最锋利的问题。延迟接受是 strategy-proof 的吗？对主动提议侧：是——谎报偏好不能让任何提议者获得比自己诚实报告时更稳定的匹配。但对被动接收侧：不一定。某家医院如果少报容量或少报对某个申请人的偏好，可能诱导算法分给自己更好的结果。这说明``提议权''不仅仅是公平问题——它直接关系到参与者有没有说谎的动力。

\textbf{不可接受的对象。}真实场景中，申请人可能不愿去某些机构，机构也可能拒绝某些申请人。偏好列表在某个截止点之后全是``不可接受''。算法只需让参与者永不向标注为不可接受的对象提议，且接收方可以在无人可接受时宁愿空缺。

\textbf{稳定匹配可能极不公平。}同一侧的两个人，一个拿到第一志愿，另一个拿到倒数第一志愿——在没有阻塞对的意义下这仍旧是``稳定''的。稳定只是阻止了局部背叛，不承诺任何全局公平性。真实市场设计中的机制（比如医院--住院医匹配系统 NRMP）额外叠加了公平约束、容量约束和防策略操纵的设计——这些已超出基础定理的范围。

画一条线：稳定匹配定理解决的是``怎样让匹配不会自己散架''这个问题。它漂亮、简单、证明只有两段，但它回答的也只是这一个问题。当你希望匹配既稳定又公平又防操纵时，你已经在做市场设计，不再只是做图论了。

\subsection{延伸与来源}

\begin{itemize}
\tightlist
\item D. Gale and L. S. Shapley, ``College Admissions and the Stability of Marriage,'' \emph{The American Mathematical Monthly}, 1962.
\item UC Berkeley CS 70: Discrete Mathematics and Probability Theory, Stable Matching 讲义。
\item MIT 6.042J / 18.062J: Mathematics for Computer Science, Matching 章节。
\end{itemize}
